\chapter{The Buy Side}\label{ch:m1:the-buy-side}

Four hundred dollars leave a teacher's payslip on the last day of the month
and arrive at her pension fund. The fund's board has decided that sixty
percent of its money tracks a world equity index, that a quarter is given to
managers paid to beat one, and that a twentieth goes to \omterm{def:m1:what-a-trading-firm-does:hedge-fund}{hedge funds}. By the
time her four hundred dollars are invested they have been split between a
machine that must buy whatever the index holds, a stock picker who is
measured every quarter, and a partnership that keeps a fifth of what it
earns. None of the three decided to trade today because of a view on prices.
This chapter is about the people on the other side of the dealer's phone:
who they are, what they are paid for, and --- the question a trader cares
about --- \emph{why they trade}.

\section{Asset owners and asset managers}

\begin{definition}[Buy side]\label{def:m1:the-buy-side:buy-side}
The \emph{buy side}\index{buy side} is the set of institutions that invest
money for a return and buy trading services to do so: \emph{asset owners}
(pension funds, insurers, sovereign funds, endowments, households), who bear
the final result, and the \emph{\omterm{def:m1:what-a-trading-firm-does:asset-manager}{asset managers}} and funds they hire.
\end{definition}

\begin{definition}[Mandate]\label{def:m1:the-buy-side:mandate}
A \emph{mandate}\index{mandate} is the contract by which an asset owner
hands a portfolio to a manager. It fixes the investable universe, the
\omterm{def:m1:the-buy-side:benchmark}{benchmark}, the risk limits, the fee, and what the manager may not do
(leverage, short sales, derivatives, concentration).
\end{definition}

\begin{remark}[Public funds are mandated by law]\label{rem:m1:the-buy-side:law}
A fund sold to the general public carries a \omterm{def:m1:the-buy-side:mandate}{mandate} written by the
legislator. A US fund that calls itself \emph{diversified} under the
Investment Company Act of 1940 must keep 75\% of its assets such that no
issuer exceeds 5\% of the fund or 10\% of that issuer's voting securities. A
European UCITS fund may put at most 10\% in one issuer, and its positions
above 5\% may not total more than 40\%. These rules, not opinions about
prices, explain a good part of why the largest funds hold what they hold ---
and why they must sell a stock that has simply risen too much.
\end{remark}

\begin{omfigure}
\begin{tikzpicture}[font=\small,
  box/.style={draw=omDef,rounded corners=2pt,minimum width=2.4cm,minimum height=0.9cm,align=center,fill=omDef!4},
  flow/.style={-{Stealth[length=2mm]},thick,omDef}]
\node[box] (sv) at (0,0) {Saver};
\node[box] (ow) at (3.5,0) {Asset owner};
\node[box] (mg) at (7.0,0) {Asset manager\\or fund};
\node[box] (sl) at (10.5,0) {Sell side};
\node[box] (mk) at (13.5,0) {Market};
\draw[flow] (sv) -- (ow);
\draw[flow] (ow) -- node[above=11pt,font=\footnotesize]{mandate} (mg);
\draw[flow] (mg) -- node[above=11pt,font=\footnotesize]{orders} (sl);
\draw[flow] (sl) -- (mk);
\node[omThm,font=\footnotesize,align=center] at (3.5,-1.05) {administration\\costs};
\node[omThm,font=\footnotesize,align=center] at (7.0,-1.05) {management and\\performance fees};
\node[omThm,font=\footnotesize,align=center] at (10.5,-1.05) {commissions, spread,\\financing};
\end{tikzpicture}

\omcaption{The chain from a saver to the market. Each link is paid (in red), and
each adds a reason to trade that has nothing to do with the price of any
security.}
\end{omfigure}

\section{Benchmarks and tracking error}

\begin{definition}[Benchmark, active return, tracking error]\label{def:m1:the-buy-side:benchmark}
A \emph{benchmark}\index{benchmark} is the reference portfolio, usually a
published index, against which a \omterm{def:m1:the-buy-side:mandate}{mandate} is measured. If $r^P_t$ and $r^B_t$
are the returns of the portfolio and the benchmark, the \emph{active return}
is $r^A_t = r^P_t - r^B_t$, and the
\emph{tracking error}\index{tracking error} is the annualised standard
deviation of $r^A$.
\end{definition}

\begin{definition}[Passive and active management]\label{def:m1:the-buy-side:passive}
\emph{Passive management}\index{passive management} aims at a \omterm{def:m1:the-buy-side:benchmark}{tracking error}
of zero: the manager holds the \omterm{def:m1:the-buy-side:benchmark}{benchmark} and is paid a few basis points for
doing it cheaply. \emph{Active management}\index{active management} accepts
a \omterm{def:m1:the-buy-side:benchmark}{tracking error} in the hope of a positive mean active return.
\end{definition}

\begin{proposition}[Tracking error from active weights]\label{prop:m1:the-buy-side:te}
Let $w^P, w^B$ be the weight vectors of portfolio and \omterm{def:m1:the-buy-side:benchmark}{benchmark} and $\Sigma$
the annualised covariance matrix of the assets' returns. With active weights
$a = w^P - w^B$ (which sum to zero), the \omterm{def:m1:the-buy-side:benchmark}{tracking error} is
$\mathrm{TE} = \sqrt{a^\top \Sigma\, a}$.
\end{proposition}
\begin{proof}
$r^A = (w^P - w^B)^\top r = a^\top r$, whose variance is $a^\top\Sigma a$.
\end{proof}

\begin{example}[Twenty overweight bets]\label{ex:m1:the-buy-side:bets}
A manager overweights 20 stocks by 1\% each and underweights 20 others by
1\% each. If stock-specific returns are independent with volatility 25\%
and the common factors cancel, $\mathrm{TE} = \sqrt{40 \times (0.01 \times
0.25)^2} = 1.6\%$. A \omterm{def:m1:the-buy-side:mandate}{mandate} that caps the \omterm{def:m1:the-buy-side:benchmark}{tracking error} at 3\% leaves room
for little more than this: an ``active'' portfolio is mostly the index.
\end{example}

\begin{definition}[Information ratio]\label{def:m1:the-buy-side:ir}
The \emph{information ratio}\index{information ratio} of a \omterm{def:m1:the-buy-side:mandate}{mandate} is the
mean annual active return divided by the \omterm{def:m1:the-buy-side:benchmark}{tracking error}:
$\mathrm{IR} = \E[r^A]/\mathrm{TE}$.
\end{definition}

\begin{proposition}[How long it takes to prove skill]\label{prop:m1:the-buy-side:years}
With independent, identically distributed active returns and a true
\omterm{def:m1:the-buy-side:ir}{information ratio} $\mathrm{IR}$, the $t$-statistic of the mean active return
after $T$ years is, in expectation, $\mathrm{IR}\sqrt{T}$. Reaching $t = 2$
takes $T = 4/\mathrm{IR}^2$ years.
\end{proposition}
\begin{proof}
Over $T$ years the sample mean of the annual active return has standard
error $\mathrm{TE}/\sqrt{T}$, so $t = \E[r^A]\sqrt{T}/\mathrm{TE} =
\mathrm{IR}\sqrt{T}$.
\end{proof}

\begin{example}[Sixteen years]\label{ex:m1:the-buy-side:sixteen}
A long-only manager with a genuinely good $\mathrm{IR}$ of 0.5 needs 16
years before its record is distinguishable from luck at the usual threshold;
with $\mathrm{IR} = 0.25$, 64 years. A market-making strategy with a daily
Sharpe ratio equivalent to $\mathrm{IR} = 8$ needs about three weeks. This one
line explains why long-horizon managers are hired on stories and fired on
noise, while short-horizon firms can be run on statistics
(\cref{fig:m1:the-buy-side:years}).
\end{example}

\begin{omfigure}
\begin{tikzpicture}
\begin{axis}[width=10.5cm,height=4.8cm,ymode=log,
  xlabel={information ratio},ylabel={years to reach $t=2$},
  tick label style={font=\small},label style={font=\small},
  xmin=0.2,xmax=3,ymin=0.4,ymax=110,grid=major,grid style={gray!25},
  ytick={0.5,1,2,5,10,20,50,100},yticklabels={0.5,1,2,5,10,20,50,100},
  table/col sep=comma]
\addplot[omDef,very thick] table[x=ir,y=years]{figdata/markets-1/03-the-buy-side/years_needed.csv};
\addplot[only marks,mark=*,omThm] coordinates {(0.5,16) (1,4) (2,1)};
\node[omThm,font=\footnotesize,anchor=west] at (axis cs:0.55,19) {16 years};
\node[omThm,font=\footnotesize,anchor=west] at (axis cs:1.05,5) {4 years};
\node[omThm,font=\footnotesize,anchor=west] at (axis cs:2.05,1.25) {1 year};
\end{axis}
\end{tikzpicture}

\omcaption{Track record needed to establish skill, $T = 4/\mathrm{IR}^2$
(logarithmic vertical axis). Data: computed by the chapter's
script.}\label{fig:m1:the-buy-side:years}
\end{omfigure}

\begin{dated}{2026-09}{How passive the market has become}\label{dat:m1:the-buy-side:passive}
At the end of 2025, index mutual funds and index exchange-traded funds held
52\% of the assets of US long-term funds, a majority for the second year, by
the fund industry association's count.
\end{dated}

\section{Hedge funds and their fees}

A \omterm{def:m1:what-a-trading-firm-does:hedge-fund}{hedge fund} (\cref{def:m1:what-a-trading-firm-does:hedge-fund}) escapes the
public-fund rules by being offered only to professional and wealthy
investors; it may borrow, sell short and use derivatives, and it is paid
differently.

\begin{definition}[Management fee, performance fee, high-water mark]\label{def:m1:the-buy-side:fees}
A \emph{management fee}\index{management fee} is a fixed annual percentage of
the assets. A \emph{performance fee}\index{performance fee} is a percentage
of the profit of the period. Under a
\emph{high-water mark}\index{high-water mark} the performance fee is charged
only on the amount by which the investor's net asset value exceeds its
highest previous fee-paying level: losses must be recovered before the
manager is paid on gains again.
\end{definition}

\begin{definition}[Commodity trading advisor]\label{def:m1:the-buy-side:cta}
A \emph{commodity trading advisor}\index{commodity trading advisor} (CTA) is,
in US law, a firm registered to advise on futures trading; in industry usage,
a fund that trades futures systematically across asset classes, most often
following trends (One Quant Book 8).
\end{definition}

\begin{method}[Computing fees under a high-water mark]\label{met:m1:the-buy-side:hwm}
Start with net asset value $N_0 = 1$ and mark $H_0 = 1$. Each year, with
gross return $g_t$, management rate $\mu$ and performance rate $\phi$:
\begin{enumerate}
\item charge the \omterm{def:m1:the-buy-side:fees}{management fee} $\mu N_{t-1}$;
\item compute $N' = N_{t-1}(1+g_t) - \mu N_{t-1}$;
\item charge the \omterm{def:m1:the-buy-side:fees}{performance fee} $\phi\,\max(0,\, N' - H_{t-1})$;
\item set $N_t = N'$ less that fee, and $H_t = \max(H_{t-1}, N_t)$.
\end{enumerate}
Conventions differ (fee on opening or average assets, quarterly
crystallisation, hurdle rates): read the offering document.
\end{method}

\begin{remark}[A performance fee is an option]\label{rem:m1:the-buy-side:option}
The payoff $\phi\max(0, N'-H)$ is that of $\phi$ call options on the fund's
value struck at the \omterm{def:m1:the-buy-side:fees}{high-water mark}, given to the manager every year for
free. An option is worth more when volatility is higher: a manager with no
skill at all earns a larger expected fee by taking more risk
(\cref{fig:m1:the-buy-side:convex}), and a manager far below its mark holds
an option so far out of the money that closing the fund and starting another
is the rational move. Investors answer with risk limits, with clawbacks, and
by demanding that the manager's own money sit in the fund.
\end{remark}

\begin{omfigure}
\begin{tikzpicture}
\begin{axis}[width=10.5cm,height=4.6cm,
  xlabel={annual volatility of the fund (\%)},ylabel={mean fee (\% a year)},
  tick label style={font=\small},label style={font=\small},
  xmin=0,xmax=30,ymin=1.5,ymax=3,grid=major,grid style={gray!25},table/col sep=comma]
\addplot[omDef,very thick,mark=*,mark size=1.3pt] table[x=vol_pct,y=fee_pct_per_year]{figdata/markets-1/03-the-buy-side/fee_vs_vol.csv};
\end{axis}
\end{tikzpicture}

\omcaption{Total fees of a ``2 and 20'' fund with a \omterm{def:m1:the-buy-side:fees}{high-water mark} and
\emph{zero} expected gross return, against the volatility it runs. Skill is
absent by construction: the rise is the value of the option. Data:
simulation by the chapter's script, 4\,000 ten-year paths per
point.}\label{fig:m1:the-buy-side:convex}
\end{omfigure}

\begin{dated}{2026-09}{The size of the hedge-fund industry}\label{dat:m1:the-buy-side:hf}
\omterm{def:m1:what-a-trading-firm-does:hedge-fund}{Hedge funds} managed about \$5.6 trillion at the end of the second quarter of
2026, after passing \$5 trillion for the first time at the end of 2025,
according to the most quoted industry database. That is under 4\% of
professionally managed assets (\cref{dat:m1:what-a-trading-firm-does:aum}),
but a far larger share of trading volume, because the capital is leveraged
and turns over many times a year.
\end{dated}

\section{Why each of them trades}

A \omterm{def:m1:what-a-trading-firm-does:market-maker}{market maker}'s first question about an order is whether the sender knows
something (\cref{def:m1:what-a-trading-firm-does:adverse-selection}). The
\omterm{def:m1:the-buy-side:buy-side}{buy side}'s motives sort themselves along that line.

\begin{example}[Motives, and what they tell the other side]\label{ex:m1:the-buy-side:motives}
\begin{center}\small
\begin{tabular}{@{}p{3.3cm}p{5.3cm}p{5.6cm}@{}}
\hline
Who & Why the order exists & What it carries \\
\hline
Index fund & cash came in or out; the index changed & no view; size and date predictable \\
Pension fund, insurer & rebalancing to target weights; hedging liabilities & no view on the stock; sells what rose \\
Active long-only & a view over months & some information, slow to matter \\
\omterm{def:m1:what-a-trading-firm-does:hedge-fund}{Hedge fund}, fundamental & a view over weeks; an event & information; urgency around news \\
\omterm{def:m1:what-a-trading-firm-does:hedge-fund}{Hedge fund}, systematic & a signal over hours to days & short-lived information; many small orders \\
Any leveraged fund & a margin call or a risk limit & no view --- but forced, large, and correlated with others' \\
Retail & savings, opinions, entertainment & almost none \\
\hline
\end{tabular}
\end{center}
The third column is why a dealer treats identical orders from different
clients differently, and why brokers sell retail flow
(\cref{ch:m1:retail-flow-and-wholesaling}).
\end{example}

\begin{remark}[Predictable flow is a strategy for someone else]\label{rem:m1:the-buy-side:predictable}
Flows that are uninformed \emph{and} predictable --- index changes, month-end
rebalancing, the daily reset of leveraged funds --- are the raw material of
the event strategies of One Quant Book 8. The \omterm{def:m1:the-buy-side:buy-side}{buy side}'s \omterm{def:m1:the-buy-side:mandate}{mandates} are
public; so, to a good approximation, are its trades.
\end{remark}

\section{Tutorial: three funds, ten years}\label{tut:m1:the-buy-side:threefunds}

\begin{tutorial}
\textbf{Goal.} Measure \omterm{def:m1:the-buy-side:benchmark}{tracking error} on simulated returns, then run a
\omterm{def:m1:what-a-trading-firm-does:hedge-fund}{hedge fund}'s fees year by year. \textbf{End state:} the chart below and the
table of the weekend problem.
\begin{tutsteps}
\item \textbf{\omterm{def:m1:the-buy-side:benchmark}{Tracking error} and \omterm{def:m1:the-buy-side:ir}{information ratio}.}
\omcode{code/markets-1/03-the-buy-side/python/buyside.py}{8}{21}{Tracking error, information ratio and the years needed to establish skill.}\label{lst:m1:the-buy-side:te}
\item \textbf{Check on simulated funds.} The figure script draws ten years
of monthly \omterm{def:m1:the-buy-side:benchmark}{benchmark} returns, an index fund (5 basis points of fee) and an
active fund built with 1\% of gross alpha, a 0.75\% fee and 4\% of
\omterm{def:m1:the-buy-side:benchmark}{tracking error}. Measured: \omterm{def:m1:the-buy-side:benchmark}{tracking errors} of 0.11\% and 3.78\%, and a
realised active return of $+0.66\%$ a year for the active fund against a
true $+0.25\%$ after fees. Ten years say almost nothing about which is
which --- as \cref{prop:m1:the-buy-side:years} predicts.
\item \textbf{Fees with a \omterm{def:m1:the-buy-side:fees}{high-water mark}.}
\omcode{code/markets-1/03-the-buy-side/python/buyside.py}{32}{51}{The method of this chapter, line for line.}\label{lst:m1:the-buy-side:hwm}
\item \textbf{Run the ten-year path} of the weekend problem. The fund earns
\omterm{def:m1:the-buy-side:fees}{performance fees} in only three years out of ten; the other seven are spent
under the mark.
\end{tutsteps}
\textbf{What to change next.} Move the \omterm{def:m1:the-buy-side:fees}{management fee} to 1\% and the
\omterm{def:m1:the-buy-side:fees}{performance fee} to 30\%: who gains? Then add a hurdle: no \omterm{def:m1:the-buy-side:fees}{performance fee} on
the first 4\% of annual return.
\end{tutorial}

\begin{omfigure}
\begin{tikzpicture}
\begin{axis}[width=11cm,height=5.2cm,
  xlabel={year},ylabel={value of 1 invested},
  tick label style={font=\small},label style={font=\small},
  xmin=0,xmax=10,ymin=0.8,ymax=2.0,grid=major,grid style={gray!25},
  legend columns=3,legend style={font=\footnotesize,at={(0.5,-0.32)},anchor=north,draw=none,fill=none,/tikz/every even column/.append style={column sep=8pt}},
  table/col sep=comma]
\addplot[omExo,thick,dashed] table[x=year,y=gross]{figdata/markets-1/03-the-buy-side/hwm_path.csv};
\addplot[omDef,very thick] table[x=year,y=net]{figdata/markets-1/03-the-buy-side/hwm_path.csv};
\addplot[omThm,thick,const plot] table[x=year,y=hwm]{figdata/markets-1/03-the-buy-side/hwm_path.csv};
\legend{gross of fees,investor's net value,high-water mark}
\end{axis}
\end{tikzpicture}

\omcaption{Ten years of a ``2 and 20'' fund. The mark rises only when the net
value makes a new high; the gap between the dashed and the solid line is
what the fees, and the growth they would have earned, cost the investor.
Data: the weekend problem's return path, computed by the chapter's script.}
\end{omfigure}

\section{Build: the fee engine}\label{bld:m1:the-buy-side:fees}

\begin{build}
\bfield{Purpose} The miniature firm will run a fund. It needs to compute
what its investors owe, per investor and per period.
\bfield{Interface} \texttt{FeeTerms(mgmt, perf, hurdle=0.0,
periods\_per\_year=1)} and \texttt{accrue(terms, state, gross\_return)}
returning a new \texttt{InvestorState(nav, high\_water)} and the two fees.
\bfield{Rules} One state per investor and subscription date: two investors
in the same fund have different marks. The \omterm{def:m1:the-buy-side:fees}{management fee} is charged on
opening value, pro rata per period. A hurdle raises the mark by the hurdle
rate each period before the \omterm{def:m1:the-buy-side:fees}{performance fee} is computed. Fees are never
negative.
\bfield{Acceptance tests} \texttt{code/firm/fees/tests/}: reproduces the
tutorial's path; no \omterm{def:m1:the-buy-side:fees}{performance fee} under the mark; with quarterly periods
and zero volatility the annual fee equals the annual rule's.
\bfield{Stretch} Post fees to the ledger of
\cref{ch:m1:what-a-trading-firm-does} as revenue of the management company.
\end{build}

\begin{omsources}
\item Investment Company Institute, \emph{2026 Investment Company Fact Book},
chapter 2.
\item HFR, \emph{Global Hedge Fund Industry Report}, second quarter 2026, and
the accompanying market commentaries.
\item US Securities and Exchange Commission staff, \emph{Report to Congress
regarding threshold limits applicable to diversified companies}, February
2022 (the 75--5--10 test of the Investment Company Act of 1940).
\item Directive 2009/65/EC (UCITS), article 52, and ESMA, \emph{Questions and
answers on the application of the UCITS Directive}.
\item R.~Grinold and R.~Kahn, \emph{Active Portfolio Management}, 2nd ed.,
McGraw-Hill, 2000, chapters 4 and 5 (information ratio, tracking error).
\item W.~Goetzmann, J.~Ingersoll and S.~Ross, ``High-water marks and hedge
fund management contracts'', \emph{Journal of Finance} 58 (2003).
\end{omsources}

\section{Exercises}

\begin{exercise}[$\star$]\label{exo:m1:the-buy-side:1}
A fund returned 9.1\%, 3.4\%, $-6.2\%$, 12.0\% in four years when its
\omterm{def:m1:the-buy-side:benchmark}{benchmark} returned 8.0\%, 4.0\%, $-7.0\%$, 10.5\%. Compute the active
returns, their mean, the \omterm{def:m1:the-buy-side:benchmark}{tracking error} (sample standard deviation) and the
\omterm{def:m1:the-buy-side:ir}{information ratio}.
\end{exercise}

\begin{exercise}[$\star$]\label{exo:m1:the-buy-side:2}
How many years of record does it take to reach $t = 2$ with an \omterm{def:m1:the-buy-side:ir}{information
ratio} of 0.4? of 1.5? How many \emph{trading days} with an annualised ratio
of 6?
\end{exercise}

\begin{exercise}[$\star$]\label{exo:m1:the-buy-side:3}
A \$50 billion index fund charges 4 basis points; a \$2 billion \omterm{def:m1:what-a-trading-firm-does:hedge-fund}{hedge fund}
charges 2 and 20 and earns 10\% gross. Which manager earns more this year?
\end{exercise}

\begin{exercise}[$\star\star$]\label{exo:m1:the-buy-side:4}
A UCITS fund holds one stock at 9.5\%, three at 6\% each and two at 5.5\%
each; every other position is below 5\%. Does it respect the 5/10/40 rule?
The 9.5\% stock rises 30\% while the rest of the fund is flat: what is its
new weight, and what must the manager do?
\end{exercise}

\begin{exercise}[$\star\star$]\label{exo:m1:the-buy-side:5}
A manager runs 50 overweights and 50 underweights of 0.5\% each;
stock-specific volatility is 30\% and specific returns are independent.
Compute the \omterm{def:m1:the-buy-side:benchmark}{tracking error}. It wants 3\%: what active weight per position
achieves it?
\end{exercise}

\begin{exercise}[$\star\star$]\label{exo:m1:the-buy-side:6}
A fund under a ``2 and 20'' contract with a \omterm{def:m1:the-buy-side:fees}{high-water mark} returns
$+30\%$, $-20\%$, $+10\%$ gross. Apply \cref{met:m1:the-buy-side:hwm} and
give, for each year, the two fees and the closing net value.
\end{exercise}

\begin{exercise}[$\star\star\star$]\label{exo:m1:the-buy-side:7}
\emph{Coding.} With \texttt{run\_fees}, compare ``2 and 20'' with ``1 and
30'' on the weekend problem's ten-year path. Report total fees under each
contract and the investor's final value. Which contract does an investor who
believes in the manager prefer, and why might the manager prefer the other?
\end{exercise}

\begin{exercise}[$\star\star\star$]\label{exo:m1:the-buy-side:8}
\emph{Find the flaw.} A consultant ranks 400 active funds by their five-year
active return, selects the top 20, and reports that ``the selected managers
have a demonstrated \omterm{def:m1:the-buy-side:ir}{information ratio} above 1''. The typical fund in the
sample has a \omterm{def:m1:the-buy-side:benchmark}{tracking error} of 4\%. Assume that \emph{no} manager has any
skill. Estimate the active return of the twentieth-best fund and its
apparent \omterm{def:m1:the-buy-side:ir}{information ratio}, and state the flaw.
\end{exercise}

\section{Problem: Two and Twenty}

\begin{problem}[{Weekend problem --- ten years in a hedge fund}]\label{pb:m1:the-buy-side:1}
An endowment invests \$100 million in a \omterm{def:m1:what-a-trading-firm-does:hedge-fund}{hedge fund} charging a 2\% \omterm{def:m1:the-buy-side:fees}{management
fee} on opening value and a 20\% \omterm{def:m1:the-buy-side:fees}{performance fee} with a \omterm{def:m1:the-buy-side:fees}{high-water mark}
(\cref{met:m1:the-buy-side:hwm}). Over ten years the fund's gross returns
are
\[
+20\%,\ -15\%,\ +10\%,\ +25\%,\ -5\%,\ +8\%,\ +30\%,\ -20\%,\ +15\%,\ +12\%.
\]

\medskip\noindent\textbf{Part I --- The first three years.}
\begin{enumerate}
\item Year 1: compute the \omterm{def:m1:the-buy-side:fees}{management fee}, the \omterm{def:m1:the-buy-side:fees}{performance fee}, the closing
net value and the new mark.
\item Year 2: the same. Why is there no \omterm{def:m1:the-buy-side:fees}{performance fee}?
\item Year 3: the fund gains 10\% gross. Is a \omterm{def:m1:the-buy-side:fees}{performance fee} due?
\item After three years, what has the investor earned, and what has the
manager been paid?
\end{enumerate}

\medskip\noindent\textbf{Part II --- The whole path} (use the tutorial's code
or a spreadsheet).
\begin{enumerate}[resume]
\item In which years is a \omterm{def:m1:the-buy-side:fees}{performance fee} paid?
\item Give the total \omterm{def:m1:the-buy-side:fees}{management fees} and the total \omterm{def:m1:the-buy-side:fees}{performance fees} over
the ten years, in millions.
\item Give the investor's final value and annualised net return.
\item Give the gross final value of \$100 million and the annualised gross
return.
\item Fees paid plus the investor's profit do not add up to the gross
profit. Compute the difference and explain it.
\end{enumerate}

\medskip\noindent\textbf{Part III --- Who got what.}
\begin{enumerate}[resume]
\item What fraction of the gross profit was paid to the manager? What
fraction did the investor keep?
\item An index fund charging 5 basis points returned 6\% a year gross over
the same decade. Compare the investor's outcome.
\item The \omterm{def:m1:what-a-trading-firm-does:hedge-fund}{hedge fund}'s returns were uncorrelated with equities. Give one
reason the endowment might still be satisfied.
\item The fund's annual gross returns have a standard deviation of 17\%.
Estimate its gross Sharpe ratio taking the mean gross return and a zero
risk-free rate. How many years would establish it at $t = 2$?
\end{enumerate}

\medskip\noindent\textbf{Part IV --- Incentives.}
\begin{enumerate}[resume]
\item At the end of year 8 the net value is far below the mark. By how much
must the fund rise, net of the \omterm{def:m1:the-buy-side:fees}{management fee}, before the manager earns a
\omterm{def:m1:the-buy-side:fees}{performance fee} again?
\item Explain why the manager might then (a) raise risk or (b) close the
fund, and who loses in each case.
\item A second investor subscribes at the start of year 9. What is
\emph{its} \omterm{def:m1:the-buy-side:fees}{high-water mark}, and what \omterm{def:m1:the-buy-side:fees}{performance fee} does it pay in year 9?
\item Why does this make per-investor accounting
(\cref{bld:m1:the-buy-side:fees}) unavoidable?
\item Using \cref{fig:m1:the-buy-side:convex}, estimate the annual fee an
unskilled manager collects at 20\% volatility, and the excess over the
\omterm{def:m1:the-buy-side:fees}{management fee} alone.
\item State the \emph{named result}: the manager's share of the gross profit
over the ten years, in percent.
\item In two sentences, say what an investor should negotiate first: the
2, the 20, or something else.
\end{enumerate}
\end{problem}

\section{Interview questions}

\begin{interviewq}[$\star$ \iqroles{trader, researcher, bank}]\label{iq:m1:the-buy-side:1}
What is \omterm{def:m1:the-buy-side:benchmark}{tracking error}, and why would a manager who is confident in a stock
still refuse to make it 15\% of the portfolio?
\end{interviewq}

\begin{interviewq}[$\star$ \iqroles{researcher, mle}]\label{iq:m1:the-buy-side:2}
A strategy has a Sharpe ratio of 1. How long a backtest do you need before
you believe it is not zero?
\end{interviewq}

\begin{interviewq}[$\star\star$ \iqroles{trader, researcher}]\label{iq:m1:the-buy-side:3}
You receive two identical orders to sell 500\,000 shares, one from an index
fund and one from a fundamental \omterm{def:m1:what-a-trading-firm-does:hedge-fund}{hedge fund}. Do you price them the same?
\end{interviewq}

\begin{interviewq}[$\star\star$ \iqroles{researcher, bank}]\label{iq:m1:the-buy-side:4}
Explain the \omterm{def:m1:the-buy-side:fees}{high-water mark} to a non-specialist, then explain what it does
to the manager's appetite for risk.
\end{interviewq}

\begin{interviewq}[$\star\star$ \iqroles{researcher, trader}]\label{iq:m1:the-buy-side:5}
Index funds hold more than half of US fund assets. Name two predictable
trading flows this creates and who might profit from each.
\end{interviewq}

\begin{interviewq}[$\star\star\star$ \iqroles{researcher, mle}]\label{iq:m1:the-buy-side:6}
Out of 1\,000 strategies with no skill, how good does the best one look
after three years? Give an order of magnitude for its Sharpe ratio.
\end{interviewq}
